\documentclass[11pt,a4paper]{article}
\usepackage[margin=25mm]{geometry}
\usepackage{amsmath,amssymb,amsthm,hyperref}
\newtheorem{lemma}{Lemma}
\newtheorem{proposition}[lemma]{Proposition}
\newtheorem{corollary}[lemma]{Corollary}
\title{A shortest-odd-cycle reduction for biaffine extremality}
\author{Human Compression project: bounded L025 proof}
\date{}
\begin{document}\maketitle
\begin{abstract}
The fixed preprint's Theorem 4.3 uses a lower bound for bipartite
competitors to make an unrestricted extremality assertion. We prove an
elementary odd-cycle bound which settles every even prime-power degree
at least eight. Together with the retained degree-four computation and
odd-degree parity, this supplies the missing unrestricted extremality
step for every prime power $q\geq3$. This note proves the claimed result
with its proper range; it does not declare the source theorem false.
\end{abstract}

\section{A uniform elementary bound}
All graphs are finite and simple. The odd girth of a nonbipartite graph
is the length of its shortest odd cycle.
\begin{lemma}\label{path}
Let $C$ be a shortest odd cycle, of length $\ell$, and let a path of
length $r$ have distinct endpoints on $C$ and all its internal vertices
outside $C$. One of the two arcs of $C$ between its endpoints has length
at most $r$ and parity equal to $r$.
\end{lemma}
\begin{proof}
Exactly one of the two arcs has parity different from $r$. Its union
with the path is a simple odd cycle, so its length is at least $\ell$.
The complementary arc consequently has length at most $r$. Since
$\ell$ is odd, this complementary arc has the same parity as $r$.
\end{proof}

\begin{proposition}\label{bound}
A nonbipartite $k$-regular graph of girth at least six and odd girth
$\ell$ has
\[
 |V|\ \geq\ \ell(k-1)+
       \left\lceil\frac{\ell(k-2)^2}3\right\rceil
 =\left\lceil\frac{\ell(k^2-k+1)}3\right\rceil .             \tag{1}
\]
In particular $\ell\geq7$.
\end{proposition}
\begin{proof}
Index a shortest odd cycle by $c_i$, $i\in\mathbb Z/\ell\mathbb Z$.
It has no chord: a chord and one of its two arcs would give a shorter
odd cycle. No outside vertex can have two neighbors on $C$. Indeed,
Lemma~\ref{path} for a path of length two supplies a two-edge arc,
giving a quadrilateral. Thus the first shell $S$, consisting of all
neighbors of $C$ outside $C$, is the disjoint union of sets $S_i$,
each of size $k-2$, with every element of $S_i$ adjacent to $c_i$ alone
on the cycle.

Suppose $a\in S_i$ and $b\in S_j$ are adjacent. Equality $i=j$ would
give a triangle. Otherwise the path $c_i,a,b,c_j$ and
Lemma~\ref{path} give an arc of odd length at most three. Length one
would give a quadrilateral. Therefore $j-i=\pm3$ modulo $\ell$.
A vertex $a\in S_i$ has at most one neighbor in either $S_{i+3}$ or
$S_{i-3}$: two in the same set would form a quadrilateral with that
set's cycle vertex. In fact it cannot have one in each. If these were
$x\in S_{i-3}$ and $z\in S_{i+3}$, the path
$c_{i-3},x,a,z,c_{i+3}$ would have length four and have all internal
vertices outside $C$. The complementary cycle arc, avoiding the six
successive forward edges from $c_{i-3}$ to $c_{i+3}$, has length
$\ell-6$, which is positive and odd. Their union is therefore a simple
odd cycle of length $\ell-2$, a contradiction. Thus the graph induced
on $S$ is a matching. Each vertex of $S$ has at most one neighbor in
$S$, one on $C$, and at least $k-2$ in
$R=V\setminus(C\cup S)$. Hence there are at least
$\ell(k-2)^2$ edges from $S$ to $R$.

For $x\in R$, two neighbors in a single $S_i$ would likewise give a
quadrilateral. If $x$ has neighbors in $S_i$ and $S_j$, the external
path of length four between $c_i,c_j$ shows that
$j-i\in\{\pm2,\pm4\}$. Thus the set of indices of its neighbors in
$S$ is a clique in the graph joining indices at differences
$\pm2,\pm4$. Multiplication by the inverse of two modulo the odd
integer $\ell$ identifies this graph with the square of the
$\ell$-cycle. Its clique number is three for $\ell\geq7$.
For clarity, in a clique containing $0$, the other indices lie in
$\{-2,-1,1,2\}$; the pairs $(-2,1),(-1,2),(-2,2)$ are not edges,
so at most two of those four indices can be chosen. The indices
$0,1,2$ exhibit a clique of size three.

Each vertex of $R$ therefore receives at most three edges from $S$.
Adding $|C|=\ell$, $|S|=\ell(k-2)$ and the resulting lower bound
on $|R|$ proves (1). The argument applies to an odd-cycle component
if the graph is disconnected; additional vertices only strengthen it.
\end{proof}

\section{Closing the source extremality argument}
The edge-rooted Moore-tree argument gives, for a girth-six
$\operatorname{egr}(v,k,6,\lambda)$,
\[
 v\geq2(k^2-k+1)+
 \left\lceil\frac{2((k-1)^3-\lambda)}k\right\rceil .       \tag{2}
\]
For bipartite graphs the last term can be replaced by twice the
ceiling with numerator $(k-1)^3-\lambda$. These bounds, and the
biaffine construction and its cycle count, are the exact inputs
already written in the parent L025 candidate. Substitution of
$k=q$, $\lambda=(q-1)^2(q-2)$ in (2) gives $v\geq2q^2-1$;
the bipartite bound gives $v\geq2q^2$. The biaffine graph attains the
latter order. Odd $q$ forces even $v$ by the handshaking identity.
Here the bounds are the corrected Moore-tree formulas, not a literal
transcription of the printed even-girth exponents in Theorems 4.1--4.2.

\begin{corollary}
For every even prime power $q\geq8$, no
$\operatorname{egr}(2q^2-1,q,6,(q-1)^2(q-2))$ exists.
Consequently the biaffine graph is unrestrictedly extremal for
those $q$, as well as for every odd prime power $q\geq3$.
\end{corollary}
\begin{proof}
A putative graph at the displayed odd order cannot be bipartite by
the bipartite lower bound. Proposition~\ref{bound} gives
$3v\geq7(q^2-q+1)$. But
\[
 7(q^2-q+1)-3(2q^2-1)=(q-5)(q-2)>0
\]
for $q\geq8$, a contradiction. The odd-degree conclusion was
established immediately above.
\end{proof}

The retained exhaustive computation of Goedgebeur and Jooken,
Table 2, row $(k,g,\lambda)=(4,6,18)$, has lower and upper bounds
both equal to $32$. Thus it excludes the remaining order $31$ at
$q=4$. This is an external finite completeness premise, not a new
execution or consequence of Proposition~\ref{bound}.
Since an even prime power is a power of two, this is the only additional
case needed. Consequently, for every prime power $q\geq3$,
\[
 n\bigl(q,6,(q-1)^2(q-2)\bigr)=2q^2.
\]
In particular the possible order $127$ at $q=8$ is excluded without a
new search: every nonbipartite $8$-regular graph of girth at least six
has at least $7(64-8+1)/3=133$ vertices. The uniform shell argument
alone gives only $31$ at $q=4$, so it does not replace the stated
external finite premise. The source's range must retain $q\geq3$,
since the biaffine graph at $q=2$ has girth eight rather than six.

\section{Scope and attribution}
This is a bounded elementary resolution of one principal
recovery gap in the fixed seventeen-page preprint, physical page 13,
Theorem 4.3 and its comparison with Theorems 4.1--4.2. It does not
revisit the accepted Suzuki calculation, supply a new Lean certificate,
add an original, or assign a numerical factor. The source's use of a
bipartite comparison bound alone is insufficient for its unrestricted
claim; that observation is a proof gap, not a theorem counterexample.
Rewrite\_termination derived the initial shell bound and the cycle-square
clique argument. Root supplied the decisive improvement that the first shell
is a matching, and rewrite\_termination independently checked and incorporated
that argument. The written shell proof is uniform and has no finite-program
premise; the degree-four case retains the named external exhaustive result. Its small optional checks are diagnostics only. No priority or
novelty claim for the girth-pair bound is made.

\begin{thebibliography}{2}
\bibitem{AL} G. Araujo-Pardo and D. Leemans,
\emph{Edge-girth-regular graphs arising from biaffine planes and Suzuki
groups}, arXiv:2108.06636v1, 15 August 2021, fixed seventeen-page edition.
\bibitem{GJ} J. Goedgebeur and J. Jooken,
\emph{Exhaustive generation of edge-girth-regular graphs},
arXiv:2401.08271v2, 25 June 2024, Table 2.
\end{thebibliography}
\end{document}
